\documentclass[12pt,a4paper]{article}
\usepackage[utf8]{inputenc}
\usepackage[T1,T2A]{fontenc}
\usepackage[ukrainian]{babel}
\usepackage{amsmath, amssymb, geometry}
\geometry{top=2cm, bottom=2cm, left=2.5cm, right=1.5cm}
\usepackage[hidelinks]{hyperref}
\usepackage{framed} % Пакет для створення блоків з лінією зліва

% Перевизначаємо leftbar для зміни товщини лінії
\renewenvironment{leftbar}{%
	\def\FrameCommand{\vrule width 0.5pt \hspace{15pt}}%
	\MakeFramed {\advance\hsize-\width \FrameRestore}}%
{\endMakeFramed}

\begin{document}
	
	\title{Геометричні та алгебраїчні основи нейронних мереж}
	\author{}
	\date{}
	\maketitle
	
	\section{Вступ: Афінні перетворення та деформація простору}
	
	Процес розпізнавання та обробки сигналів у нейронних мережах зводиться до серії тензорних згорток і матричних множень. Кожен прихований шар мережі виконує композицію афінного перетворення та нелінійного відображення, що математично записується як:
	$$y = \sigma(Wx + b)$$
	де $W$ — матриця ваг (лінійного оператора), $x$ — вектор вхідних ознак з попереднього шару, $b$ — вектор зміщення, а $\sigma$ — нелінійна функція активації.
	
	\subsection{Фундаментальна роль нелінійності та теорема Колмогорова-Арнольда}
	Присутність оператора $\sigma$ є критичною. Без нелінійності суперпозиція будь-якої кількості шарів математично колапсує в одне-єдине еквівалентне афінне перетворення, оскільки композиція лінійних відображень завжди є лінійним відображенням. Афінне перетворення $Wx + b$ виконує лише базові геометричні операції у векторному просторі (обертання, масштабування, відображення та трансляцію). Усі гіперплощини залишаються плоскими, а паралельні лінії — паралельними.
	
	З точки зору топології, нелінійна функція деформує многовид даних, розбиваючи векторний простір і генеруючи кусково-лінійні або криволінійні границі. Теоретичним фундаментом для такої архітектури виступає \textbf{теорема Колмогорова-Арнольда} (про суперпозицію). Вона стверджує, що будь-яку неперервну функцію $n$ змінних можна подати як суперпозицію неперервних функцій однієї змінної та операції додавання:
	$$f(x_1, \dots, x_n) = \sum_{q=0}^{2n} \Phi_q \left( \sum_{p=1}^n \phi_{q,p}(x_p) \right)$$
	
	\subsubsection*{Модернізація теореми: від класичного результату до сучасних архітектур}
	
	Класичний запис теореми Колмогорова був згодом вдосконалений іншими математиками, що сформувало прямий алгебраїчний міст до сучасних архітектур штучних нейронних мереж\footnote{Детальніше: \url{https://habr.com/ru/articles/856776/}}:
	
	\begin{itemize}
		\item \textbf{Теорема Лоренца (універсальна зовнішня функція):} 
		Математик Джордж Лоренц (George Lorentz) довів, що зовнішні функції $\Phi_q$ можна звести до \textit{однієї} неперервної функції $\Phi$, яка є спільною для всіх доданків. 
		\newline \textbf{Зв'язок з нейромережами:} Це відкриття стало строгим математичним обґрунтуванням того, чому в сучасних багатошарових перцептронах (MLP) усі вузли одного прихованого шару можуть використовувати \textit{однакову} універсальну функцію активації (наприклад, ReLU, Sigmoid або Tanh). Мережі не потрібно підбирати унікальну топологічну деформацію для кожного нейрона — достатньо однієї глобальної функції.
		
		\item \textbf{Теорема Шпрехера (афінні зсуви внутрішніх функцій):} 
		Девід Шпрехер (David Sprecher) пішов ще далі, довівши, що набір внутрішніх функцій $\phi_{q,p}$ можна замінити на одну універсальну монотонну функцію $\psi$, застосовуючи до її аргументів лише константи масштабування та зсуву. Рівняння набуло вигляду:
		$$f(x_1, \dots, x_n) = \sum_{q=0}^{2n} \Phi_q \left( \sum_{p=1}^n \lambda_p \psi(x_p + \epsilon_q) \right)$$
		\textbf{Зв'язок з нейромережами:} У термінах лінійної алгебри, масштабування $\lambda_p$ та зсув $\epsilon_q$ є точним еквівалентом \textit{вагових коефіцієнтів} (weights, матриця $W$) та \textit{векторів зміщення} (biases, вектор $b$). Шпрехер довів, що складність багатовимірної апроксимації можна перенести з пошуку складних функцій на простий пошук афінних параметрів масштабу і зсуву навколо однієї базової кривої $\psi$.
	\end{itemize}
	
	\subsubsection*{Аналогія з архітектурою нейронних мереж (MLP проти KAN)}
	
	Модернізація цієї теореми — зокрема, у роботах Хехт-Нільсена (Hecht-Nielsen, 1987) та в сучасній архітектурі мереж Колмогорова-Арнольда (KAN\footnote{Ziming Liu, Yixuan Wang, Sachin Vaidya, Fabian Ruehle, James Halverson, Marin Soljačić, Thomas Y. Hou, Max Tegmark. ``KAN: Kolmogorov-Arnold Networks''. 2024. LaTeX-джерело: \url{https://arxiv.org/src/2404.19756}. PDF: \url{https://arxiv.org/pdf/2404.19756}.}) — дозволяє розглядати наведену формулу не просто як абстрактну тотожність, а як точний математичний опис ідеальної двошарової нейронної мережі прямого поширення. Втім, між класичними багатошаровими перцептронами (MLP) та мережами Колмогорова-Арнольда (KAN) існує фундаментальна топологічна різниця в тому, як саме кожна архітектура реалізує елементи цього рівняння:
	
	\begin{itemize}
		\item \textbf{Вхідний вектор $x_p$}: Відповідає $n$ координатам вхідних даних (нейронам першого шару), де кожна координата обробляється незалежно і розглядається як окремий вимір простору.
		
		\item \textbf{Внутрішні функції $\phi_{q,p}(x_p)$}: Згідно з теоремою, це нелінійні одновимірні деформації простору, що застосовуються \textit{до} агрегації сигналу.
		\begin{itemize}
			\item У мережах \textbf{KAN} це реалізується буквально: на кожному синапсі (ребрі) обчислювального графа розміщується параметризована нелінійна функція (B-сплайн).
			\item У класичних \textbf{MLP} цей крок математично деградує до звичайного афінного відображення (скалярне множення на вагу $w_{q,p} \cdot x_p$), повністю втрачаючи нелінійність на етапі передачі сигналу між шарами.
		\end{itemize}
		
		\item \textbf{Внутрішня сума $\sum_{p=1}^n$}: Відповідає лінійній агрегації сигналів у вузлах прихованого шару. Геометрично це означає проектування $n$-вимірного вхідного простору у простір вищої розмірності $\mathbb{R}^{2n+1}$ (оскільки індекс $q$ змінюється від $0$ до $2n$). У мережах KAN вузли виконують \textit{виключно} цю операцію підсумовування, без жодної додаткової нелінійності.
		
		\item \textbf{Зовнішні функції $\Phi_q$}: 
		\begin{itemize}
			\item У мережах \textbf{KAN} це другий шар нелінійних функцій (сплайнів), розміщених на ребрах між прихованим та вихідним шаром.
			\item У класичних \textbf{MLP} саме тут застосовується функція активації (наприклад, ReLU, Sigmoid чи Tanh) всередині вузла \textit{після} агрегації. При цьому, завдяки теоремі Лоренца, замість $2n+1$ унікальних функцій $\Phi_q$ MLP використовує єдину універсальну функцію $\sigma$ для всіх нейронів шару.
		\end{itemize}
		
		\item \textbf{Зовнішня сума $\sum_{q=0}^{2n}$}: Фінальна лінійна комбінація, яка колапсує багатовимірний простір $\mathbb{R}^{2n+1}$ у скалярний результат $f(x)$ на вихідному нейроні.
	\end{itemize}
	
	У класичних багатошарових перцептронах (MLP) лінійні (афінні) перетворення застосовуються на ребрах (ваги матриці $W$), а нелінійні універсальні функції активації $\Phi$ (за Лоренцом) розташовуються у вузлах. Сучасна архітектура KAN концептуально «перевертає» цю геометрію: спираючись на оригінальну формулу Колмогорова-Арнольда, вона відмовляється від фіксованих функцій активації у вузлах. Натомість усі нелінійні функції $\phi_{q,p}$ (які реалізуються як параметризовані B-сплайни з навчальними коефіцієнтами) розташовуються \textit{безпосередньо на ребрах} графа обчислень. Вузли у KAN виконують лише операцію додавання (лінійну агрегацію).
	
	З геометричної точки зору, замість того, щоб проектувати простір гіперплощинами та деформувати його у вузлах (як в MLP), KAN деформує кожен одновимірний базисний напрямок безпосередньо під час передачі сигналу по тензорному ребру.
	
	\begin{leftbar}
		\itshape
		\textbf{Поглиблення: Алгебраїчна реалізація через B-сплайни} \\
		У мережах KAN кожна одновимірна функція на ребрі $\phi(x)$ не є фіксованою, а формується як лінійна комбінація елементів локального базису (B-сплайнів):
		$$\phi(x) = \sum_{i=1}^k c_i B_i(x)$$
		де $B_i(x)$ — базисні функції сплайна заданого степеня, а $c_i$ — параметри (ваги), що оптимізуються. З точки зору лінійної алгебри, це елегантний перехід: задача пошуку довільної нелінійної топологічної деформації у нескінченновимірному функціональному просторі зводиться до пошуку скінченного вектора коефіцієнтів $\mathbf{c} \in \mathbb{R}^k$ у лінійному підпросторі, натягнутому на базис $\{B_i\}$.
	\end{leftbar}
	
	Ключовий внесок цієї роботи полягає в узагальненні класичного зображення Колмогорова-Арнольда — яке відповідає лише дворівневій KAN-мережі форми $[n,\,2n+1,\,1]$ — на довільну ширину та глибину: подібно до того, як шари MLP накладаються один на одний, шари KAN також складаються у глибші мережі. Формально це записується як композиція шар-матриць одновимірних функцій $\Phi_l$:
	$$\text{KAN}(x) = (\Phi_{L-1} \circ \Phi_{L-2} \circ \cdots \circ \Phi_0)(x),$$
	що структурно контрастує з класичним MLP, який чергує афінні перетворення з фіксованою нелінійністю:
	$$\text{MLP}(x) = (W_{L-1} \circ \sigma \circ W_{L-2} \circ \sigma \circ \cdots \circ W_0)(x).$$
	
	Отже, теорема Колмогорова-Арнольда геометрично доводить, що простір будь-яких складних багатовимірних функцій можна звести до композиції одновимірних деформацій та лінійних проекцій: у контексті машинного навчання це є строгою гарантією того, що двошарова мережа з достатньою кількістю прихованих нейронів (мінімум $2n+1$) здатна апроксимувати довільну неперервну функцію на обмеженій області. Проте саме мережі KAN реалізують цю геометрію найбільш точно та природно: вони буквально розміщують передбачені теоремою одновимірні функції на ребрах графа, тоді як класичні MLP лише наближено відтворюють цю структуру за допомогою фіксованих активацій у вузлах та лінійних ваг на ребрах.
	
	\section{Класифікація нелінійних відображень}
	З точки зору лінійної алгебри, нелінійне відображення $\sigma: \mathbb{R}^n \to \mathbb{R}^m$ може мати різну структуру. Розглянемо три основні геометричні класи:
	
	\subsection{1. Ізотропні покомпонентні відображення (Diagonal maps)}
	Тут $\sigma$ — це одна й та сама скалярна функція $f: \mathbb{R} \to \mathbb{R}$, яка застосовується до кожної координати вектора $\mathbf{z} \in \mathbb{R}^n$ абсолютно незалежно:
	$$\sigma(\mathbf{z}) = \begin{pmatrix} f(z_1) \\ f(z_2) \\ \vdots \\ f(z_n) \end{pmatrix}$$
	\textbf{Геометричний зміст:} Деформація простору відбувається виключно вздовж кожного базисного вектора окремо. Взаємодії між вимірами на етапі активації немає (за змішування координат відповідає виключно лінійний оператор $W$).
	
	\subsection{2. Глобальні векторні відображення (Взаємозалежні координати)}
	Значення кожної вихідної координати залежить від усього вхідного вектора. Тут $\sigma$ є вектор-функцією, яку неможливо розділити. Прикладом є відображення Softmax:
	$$\sigma(\mathbf{z})_i = \frac{e^{z_i}}{\sum_{j=1}^{n} e^{z_j}}$$
	\textbf{Геометричний зміст:} Це гладке відображення бере весь векторний простір $\mathbb{R}^n$ і проектує його у внутрішність $(n-1)$-вимірного standard симплекса. Через спільний знаменник координати стають лінійно залежними ($\sum \sigma_i = 1$), що кардинально змінює топологію простору.
	
	\subsection{3. Анізотропні (параметричні) відображення}
	Функція діє покомпонентно, але має різну форму для різних базисних напрямків. Наприклад, у Parametric ReLU (PReLU):
	$$\sigma(\mathbf{z})_i = \max(0, z_i) + \alpha_i \min(0, z_i)$$
	де вектор $\alpha \in \mathbb{R}^n$ містить унікальні коефіцієнти розтягування/стискання для кожної $i$-тої осі.
	\textbf{Геометричний зміст:} Гіперплощини простору деформуються (згинаються) з різним ступенем жорсткості для кожного виміру. Простір зазнає анізотропної кусково-лінійної деформації.
	
	\section{Лінійна алгебра: Згортання тензорів (Tensor Contraction)}
	
	У машинному зорі вхідне зображення є тензором третього рангу $\mathcal{X} \in \mathbb{R}^{H \times W \times C}$, де $H, W$ — просторові розміри, а $C$ — кількість каналів. 
	Ядро згортки (фільтр) — це тензор четвертого рангу $\mathcal{W} \in \mathbb{R}^{K \times K \times C \times M}$, де $K$ — просторовий розмір фільтра, а $M$ — кількість вихідних каналів.
	
	Математично, для кожного пікселя вихідного тензора $\mathcal{Y} \in \mathbb{R}^{H' \times W' \times M}$, обчислюється скалярний добуток Фробеніуса (згортання за трьома індексами) між локальним патчем зображення та фільтром:
	
	$$ \mathcal{Y}_{i,j,m} = \sum_{c=1}^{C} \sum_{u=0}^{K-1} \sum_{v=0}^{K-1} \mathcal{X}_{i+u, j+v, c} \cdot \mathcal{W}_{u, v, c, m} $$
	
	Тут ми бачимо класичну афінну трансформацію (до результату також додається вектор зміщення $\mathbf{b} \in \mathbb{R}^M$), яка проектує підпростір розмірності $K \times K \times C$ у скаляр.
	
	\begin{leftbar}
		\itshape
		\textbf{Поглиблення: Тензорне числення та угода Ейнштейна} \\
		У строгому математичному апараті диференціальної геометрії використовується угода Ейнштейна про підсумовування (де сума береться за індексами, що повторюються зверху і знизу), що дозволяє записати операцію згортки без громіздких знаків суми.
		
		Нехай $X^{c}_{x, y}$ — компоненти вхідного тензора (верхній індекс $c$ — контраваріантна компонента каналу, нижні $x, y$ — коваріантні просторові координати), а $W^{m}_{c, u, v}$ — компоненти тензора ваг ($m$ — вихідний канал). Тоді операція локальної згортки набуває вигляду:
		$$ Y^{m}_{i, j} = X^{c}_{i+u, j+v} W^{m}_{c, u, v} + b^m $$
		Підсумовування за індексами $c$ та $u, v$ мається на увазі автоматично. Це чітко демонструє, що згортка є класичним згортанням тензорів за відповідними коваріантними та контраваріантними індексами.
	\end{leftbar}
	
	\subsection{Матрична репрезентація (Тоепліцеві матриці та im2col)}
	
	Апаратно обчислювати багатовимірні суми неефективно, тому тензорна згортка зводиться до одного матричного множення через ізоморфізм просторів. 
	
	Якщо розглядати одновимірну згортку, то ядро діє як лінійний оператор, матриця якого має структуру матриці Тоепліца (діагоналі містять однакові елементи). Для двовимірного випадку (зображення) це буде двічі блочно-циркулянтна матриця з Тоепліцевими блоками.
	
	На практиці застосовується операція перетворення im2col (image to column):
	\begin{enumerate}
		\item Локальні патчі тензора $\mathcal{X}$ (розміром $K \times K \times C$) розгортаються у вектори і записуються як стовпці матриці $\mathbf{X}_{col}$.
		\item Тензор ваг $\mathcal{W}$ розгортається у матрицю $\mathbf{W}_{row}$.
		\item Згортка виконується як стандартне множення матриць:
		$$ \mathbf{Y} = \mathbf{W}_{row} \mathbf{X}_{col} $$
	\end{enumerate}
	Таким чином, згортка — це лінійне відображення $\mathbb{R}^{K^2C} \to \mathbb{R}^M$, яке застосовується до кожного локального патча зображення.
	
	\subsection{Диференціальна геометрія: Згортка на многовидах та еквіваріантність}
	
	У класичному машинному зорі ми працюємо на плоскому евклідовому просторі $\mathbb{R}^2$ або дискретній ґратці $\mathbb{Z}^2$. Операція згортки є еквіваріантною (трансляційно-інваріантною) відносно групи зсувів $T(2)$. Тобто, якщо зсунути об'єкт на зображенні, карта ознак (feature map) зсунеться аналогічно: $f(T_g x) = T_g f(x)$.
	
	Проте, якщо подивитися на це через призму диференціальної геометрії, ми можемо узагальнити згортку для довільних многовидів (наприклад, сфера $S^2$ для панорамних камер) або груп Лі (наприклад, $SE(2)$ для розпізнавання об'єктів незалежно від їхнього кута повороту).
	
	У цьому контексті:
	\begin{itemize}
		\item Вхідне зображення розглядається як функція (або переріз векторного розшарування) на гладкому многовиді $M$.
		\item Згортка формалізується через інтегрування за мірою Хаара на топологічній групі $G$:
		$$ [\mathcal{W} \star \mathcal{X}](g) = \int_{G} \mathcal{W}(h^{-1}g) \mathcal{X}(h) d\mu(h) $$
		де $g, h \in G$. 
		\item У геометричному глибокому навчанні (Geometric Deep Learning) ядро згортки переміщується по многовиду за допомогою паралельного перенесення (parallel transport) уздовж геодезичних, враховуючи кривизну простору (зв'язність Леві-Чівіти), щоб порівнювати тензори у різних дотичних просторах.
	\end{itemize}
	
	\begin{leftbar}
		\itshape
		\textbf{Поглиблення: Перерізи розшарувань на групах Лі} \\
		У строгому сенсі диференціальної геометрії, вхідне зображення та карти ознак є не просто функціями, а \textbf{перерізами векторних розшарувань} (sections of vector bundles) над базовим гладким многовидом $M$. Якщо $\mathcal{E} \xrightarrow{\pi} M$ — векторне розшарування, то карта ознак $f \in \Gamma(M, \mathcal{E})$.
		
		Еквіваріантність згортки вимагає, щоб дія групи Лі $G$ на базовому многовиді $M$ коректно піднімалася (lift) до дії на просторі перерізів. Якщо $L_g$ — оператор лінійного представлення групи, то еквіваріантність формалізується як комутативність: $\mathcal{K} \star (L_g f) = L_g (\mathcal{K} \star f)$.
		
		На викривлених многовидах інтегрування за мірою Хаара $\mu$ на групі $G$ вимагає використання \textbf{паралельного перенесення} уздовж геодезичних, що дозволяє коректно порівнювати тензори у різних дотичних просторах $T_x M$ та $T_y M$.
	\end{leftbar}
	
	\section*{Джерела}
	\begin{itemize}
		\item Ziming Liu, Yixuan Wang, Sachin Vaidya, Fabian Ruehle, James Halverson, Marin Soljačić, Thomas Y. Hou, Max Tegmark. \textit{KAN: Kolmogorov-Arnold Networks}. arXiv:2404.19756 [cs.LG], 2024 (прийнято на ICLR 2025). LaTeX-джерело: \url{https://arxiv.org/src/2404.19756}. PDF: \url{https://arxiv.org/pdf/2404.19756}.
		\item Sqrudj. \textit{KAN: Kolmogorov–Arnold Networks} (повний переклад статті). Хабр, 2024. \url{https://habr.com/ru/articles/856776/}.
	\end{itemize}
	
\end{document}